%% =====================================================================
%%  Corrigés — Unit 13 « What are you doing? »
%% =====================================================================
\corrigesUnite{13}{What are you doing?}

\begin{corrige}{1}
\begin{reponses}[4]
\item \rep{working}
\item \rep{making}
\item \rep{sitting}
\item \rep{playing}
\item \rep{writing}
\item \rep{running}
\item \rep{having}
\item \rep{swimming}
\end{reponses}
\pourquoi{\en{make}, \en{write}, \en{have} perdent leur \en{-e} muet ; \en{sit}, \en{run}, \en{swim}
(une voyelle + une consonne) doublent la consonne finale.}
\piege{\en{play} se termine par un \en{y} : on ajoute simplement \en{-ing} → \en{playing}.}
\end{corrige}

\begin{corrige}{2}
\begin{reponses}[2]
\item \rep{She's talking on the phone.}
\item \rep{He's drinking a coffee.}
\item \rep{She's reading a book.}
\item \rep{He's listening to music.}
\item \rep{He's drawing a logo.}
\item \rep{She's singing.}
\end{reponses}
\retiens{\en{listen \textbf{to}} : on écoute \textbf{quelque chose} → \en{listen to music}.}
\end{corrige}

\begin{corrige}{3}
\begin{reponses}[1]
\item \rep{Is Léa talking on the phone?}
\item \rep{They aren't waiting for the bus.}
\item \rep{Is Sam wearing a suit?}
\item \rep{I'm not reading a book.}
\item \rep{Emma isn't having lunch.}
\item \rep{Are you listening to me?}
\end{reponses}
\pourquoi{pour la question, \en{be} passe devant le sujet ; pour la négation, on ajoute \en{not} après
\en{be}. Le verbe en \en{-ing} ne change pas.}
\piege{\en{I amn't} n'existe pas : on dit \en{I'm not}.}
\end{corrige}

\begin{corrige}{4}
\begin{reponses}[2]
\item \rep{wears}
\item \rep{is wearing}
\item \rep{is playing}
\item \rep{plays}
\item \rep{is sleeping}
\item \rep{don't drink}
\item \rep{are you doing}
\item \rep{Do you speak}
\end{reponses}
\pourquoi{\en{usually}, \en{every day}, \en{in the evening} signalent une habitude (présent simple) ;
\en{today}, \en{Listen!}, \en{Shh!}, \en{right now} signalent une action en cours (présent en \en{-ing}).
Phrase 8 : parler une langue est une capacité permanente, pas une action en cours.}
\end{corrige}

\begin{corrige}{5}
\begin{reponses}[3][french]
\item \rep{V}
\item \rep{F}
\item \rep{F}
\item \rep{V}
\item \rep{F}
\item \rep{V}
\end{reponses}
\pourquoi{Jack est au parc (\en{I'm at the park}) et joue au \en{soccer}, notre football. Kate ne joue
pas : elle est assise sur un banc et lit.}
\end{corrige}

\begin{corrige}{6}
\begin{reponses}[6][french]
\item \rep{b}
\item \rep{d}
\item \rep{a}
\item \rep{f}
\item \rep{c}
\item \rep{e}
\end{reponses}
\retiens{au téléphone, on se présente avec \en{This is…} : \en{This is Sam Brooks.}}
\end{corrige}

\begin{corrige}{7}
\begin{reponses}[3]
\item \rep{thing}
\item \rep{win}
\item \rep{sing}
\item \rep{rang}
\item \rep{ban}
\item \rep{sung}
\end{reponses}
\pourquoi{avec \phon{n}, le bout de la langue touche l'arrière des dents du haut ; avec \phon{ŋ}, c'est le
fond de la langue qui remonte. Texte de la piste page~\pageref{trans-13.10}.}
\end{corrige}

\begin{corrige}{8}
\begin{reponses}[5]
\item \rep{drinking}
\item \rep{reading}
\item \rep{playing}
\item \rep{working}
\item \rep{cooking}
\end{reponses}
\pourquoi{texte de la piste page~\pageref{trans-13.11}. \en{pancakes} : les crêpes épaisses du
petit-déjeuner américain.}
\end{corrige}

\begin{corrige}{9}
\begin{reponses}[4]
\item \rep{speaking}
\item \rep{speak}
\item \rep{Hold}
\item \rep{meeting}
\item \rep{leave}
\item \rep{calling}
\item \rep{This}
\item \rep{writing}
\end{reponses}
\end{corrige}

\begin{corrige}{10}
\begin{reponses}[1]
\item \rep{What are you doing? — I'm watching TV.}
\item \rep{She isn't working today.} \textit{ou} \rep{She's not working today.}
\item \rep{Listen! Sam is playing the guitar.}
\item \rep{He's wearing a suit and a tie.}
\item \rep{Hold on, please.}
\item \rep{He's in a meeting. Can I take a message?}
\item \rep{We're waiting for the bus.}
\item \rep{Are you listening to me? — Yes, I am.}
\end{reponses}
\piege{phrase 8 : pas de contraction dans une réponse courte affirmative : \en{Yes, I am}, jamais
\en{Yes, I'm}.}
\end{corrige}

\begin{corrige}{11}
Réponse modèle :
\begin{anglais}
\begin{tcolorbox}[otBase,colback=otGreenTint,colframe=otGreenTint,fontupper=\fontsize{10.4}{14}\selectfont]
Hi! I'm at a café near my office. I'm drinking a tea and I'm reading a magazine. My friend Paul is
working on his computer. Two students are talking and laughing. I'm wearing jeans, a black sweater and
my new sneakers!
\end{tcolorbox}
\end{anglais}
\retiens{\en{work on a computer} = « travailler sur un ordinateur » ; \en{laugh} = « rire ».}
\end{corrige}

\begin{corrige}{12}
Réponses modèles :
\begin{reponses}[1]
\item \rep{I'm cooking dinner.} \textit{ou} \rep{I'm reading a book.}
\item \rep{I'm at home.} \textit{ou} \rep{I'm at work.}
\item \rep{No, I'm not.} \textit{ou} \rep{Yes, I am.}
\item \rep{I'm wearing jeans and a blue sweater.}
\item \rep{Yes, I can! See you soon!} \textit{ou} \rep{Sorry, I can't. I'm busy.}
\end{reponses}
\end{corrige}

\begin{corrigetest}
\begin{reponses}[5][french]
\item \rep{b}
\item \rep{a}
\item \rep{c}
\item \rep{a}
\item \rep{b}
\item \rep{c}
\item \rep{a}
\item \rep{b}
\item \rep{c}
\item \rep{b}
\end{reponses}
\pourquoi{question 6 : \en{want} est un verbe d'état ; il reste au présent simple. Question 8 :
\en{Who are you?} est beaucoup trop direct au téléphone.}
\end{corrigetest}
